\documentclass[12pt]{article}
%^ Start of document
\usepackage{times}  %Makes text TimesNewRoman
\usepackage{setspace} %Allows you to set single or double space 
\usepackage{biblatex} %Imports biblatex package
\usepackage{graphicx} % Required for inserting images
\usepackage{authblk}
\usepackage[colorlinks=true, urlcolor=blue, linkcolor=blue]{hyperref}
\usepackage{subfig}
\usepackage{float}
\usepackage{tabularx}
\usepackage{multirow}
\usepackage[paperheight=11in,
   paperwidth=8.5in,
   top=20mm,
   bottom=20mm,
   left=40mm,
   right=40mm]{geometry}
\usepackage{moresize}
\usepackage{tikz}
\usepackage{xcolor}
\usepackage{physics}
\usepackage{amssymb}
\usepackage{mathrsfs}
\usepackage{amsmath}
\usepackage{ragged2e}
\usepackage{comment}
\usepackage{xcolor}
\addbibresource{mathmeth.bib}
\begin{document}
\singlespacing  %Sets single spacing for whole document
\title{Math Methods HW 4}

\author[1]{Zachary Tullier}
\affil[1]{Student\\Physics Department\\ University of Houston - Clear Lake \\Houston, Texas\\U.S}
\date{}
\maketitle




\section*{Textbook Question 9.3.1}
    Show that by making a change of variables to 
    \begin{gather*}
        \begin{split}
            \xi = \sqrt{c} x - \frac{b y}{\sqrt{c}}
            \\
            \eta = \frac{y}{\sqrt{c}}
        \end{split}
    \end{gather*}
    the operator $\mathcal{L}$ of Equation \ref{eq: 9.18} can be brought to the forrm 
    \[
        \mathcal{L} = (ac - b^2) \frac{\partial^2}{\partial \xi^2} + \frac{\partial^2}{\partial \eta^2}
    \]
    \begin{equation} \label{eq: 9.18}
        \mathcal{L} = \left(\frac{b + \sqrt{b^2 - ac}}{\sqrt{c}} \frac{\partial}{\partial x} + \sqrt{c} \frac{\partial}{\partial y} \right) \left(\frac{b - \sqrt{b^2 - ac}}{\sqrt{c}} \frac{\partial}{\partial x} + \sqrt{c} \frac{\partial}{\partial y} \right)
    \end{equation} 
    
    \subsection*{Solution:}
        The given change of variables is:
        \[
        \xi = c^{1/2}x - c^{-1/2}by, \quad \eta = c^{-1/2}y.
        \]
        We compute the partial derivatives of \(\xi\) and \(\eta\) with respect to \(x\) and \(y\):
        \[
        \frac{\partial \xi}{\partial x} = c^{1/2}, \quad \frac{\partial \xi}{\partial y} = -c^{-1/2}b,
        \]
        \[
        \frac{\partial \eta}{\partial x} = 0, \quad \frac{\partial \eta}{\partial y} = c^{-1/2}.
        \]
        Using the chain rule, we can express \(\frac{\partial}{\partial x}\) and \(\frac{\partial}{\partial y}\) in terms of \(\frac{\partial}{\partial \xi}\) and \(\frac{\partial}{\partial \eta}\):
        \[
        \frac{\partial}{\partial x} = \frac{\partial \xi}{\partial x} \frac{\partial}{\partial \xi} + \frac{\partial \eta}{\partial x} \frac{\partial}{\partial \eta} = c^{1/2} \frac{\partial}{\partial \xi},
        \]
        \[
        \frac{\partial}{\partial y} = \frac{\partial \xi}{\partial y} \frac{\partial}{\partial \xi} + \frac{\partial \eta}{\partial y} \frac{\partial}{\partial \eta} = -c^{-1/2}b \frac{\partial}{\partial \xi} + c^{-1/2} \frac{\partial}{\partial \eta}.
        \]
        Substitute these expressions for \(\frac{\partial}{\partial x}\) and \(\frac{\partial}{\partial y}\) into \(\mathcal{L}\):
        \[
        \mathcal{L} = \left( \frac{b + \sqrt{b^2 - ac}}{c^{1/2}} \frac{\partial}{\partial x} + c^{1/2} \frac{\partial}{\partial y} \right) \left( \frac{b - \sqrt{b^2 - ac}}{c^{1/2}} \frac{\partial}{\partial x} + c^{1/2} \frac{\partial}{\partial y} \right).
        \]
        Substituting \(\frac{\partial}{\partial x} = c^{1/2} \frac{\partial}{\partial \xi}\) and \(\frac{\partial}{\partial y} = -c^{-1/2}b \frac{\partial}{\partial \xi} + c^{-1/2} \frac{\partial}{\partial \eta}\), we expand each term.
        The first operator becomes:
        \[
        \frac{b + \sqrt{b^2 - ac}}{c^{1/2}} \cdot c^{1/2} \frac{\partial}{\partial \xi} + c^{1/2} \left(-c^{-1/2}b \frac{\partial}{\partial \xi} + c^{-1/2} \frac{\partial}{\partial \eta} \right).
        \]
        Simplify:
        \[
        \left(b + \sqrt{b^2 - ac}\right) \frac{\partial}{\partial \xi} - b \frac{\partial}{\partial \xi} + \frac{\partial}{\partial \eta}.
        \]
        Combine terms:
        \[
        \left(\sqrt{b^2 - ac}\right) \frac{\partial}{\partial \xi} + \frac{\partial}{\partial \eta}.
        \]
        The second operator becomes:
        \[
        \frac{b - \sqrt{b^2 - ac}}{c^{1/2}} \cdot c^{1/2} \frac{\partial}{\partial \xi} + c^{1/2} \left(-c^{-1/2}b \frac{\partial}{\partial \xi} + c^{-1/2} \frac{\partial}{\partial \eta} \right).
        \]
        Simplify:
        \[
        \left(b - \sqrt{b^2 - ac}\right) \frac{\partial}{\partial \xi} - b \frac{\partial}{\partial \xi} + \frac{\partial}{\partial \eta}.
        \]
        Combine terms:
        \[
        \left(-\sqrt{b^2 - ac}\right) \frac{\partial}{\partial \xi} + \frac{\partial}{\partial \eta}.
        \]
        Now, multiply the two operators:
        \[
        \mathcal{L} = \left( \left(\sqrt{b^2 - ac}\right) \frac{\partial}{\partial \xi} + \frac{\partial}{\partial \eta} \right) \left( \left(-\sqrt{b^2 - ac}\right) \frac{\partial}{\partial \xi} + \frac{\partial}{\partial \eta} \right).
        \]
        Expand:
        \[
        \mathcal{L} = (\sqrt{b^2 - ac})(-\sqrt{b^2 - ac}) \frac{\partial^2}{\partial \xi^2} + (\sqrt{b^2 - ac}) \frac{\partial}{\partial \xi} \frac{\partial}{\partial \eta}
        \]
        \[
        + (-\sqrt{b^2 - ac}) \frac{\partial}{\partial \xi} \frac{\partial}{\partial \eta} + \frac{\partial^2}{\partial \eta^2}.
        \]
        Simplify:
        \[
        \mathcal{L} = -(b^2 - ac) \frac{\partial^2}{\partial \xi^2} + \frac{\partial^2}{\partial \eta^2}.
        \]
        Factor out the coefficient of \(\frac{\partial^2}{\partial \xi^2}\):
        \[
        \mathcal{L} = (ac - b^2) \frac{\partial^2}{\partial \xi^2} + \frac{\partial^2}{\partial \eta^2}.
        \]
        We have shown that the operator \(\mathcal{L}\) can be transformed into the form:
        \[
        \mathcal{L} = (ac - b^2) \frac{\partial^2}{\partial \xi^2} + \frac{\partial^2}{\partial \eta^2},
        \]
        under the given change of variables.

\section*{Textbook Question 9.4.2}
    Show that the Helmholtz equation,
    \[
        \nabla^2 \psi + k^2 \psi = 0
    \]
    is still separable in circular cylindrical coordinates if $k^2$ is generalized to $k^2 + f(p) + \frac{1}{\rho^2} g(\varphi) + h(z)$
   
   \subsection*{Solution:}
        In cylindrical coordinates \((\rho, \phi, z)\), the Laplacian operator \(\nabla^2\) is expressed as:
        \[
        \nabla^2 \psi = \frac{1}{\rho} \frac{\partial}{\partial \rho} \left( \rho \frac{\partial \psi}{\partial \rho} \right) + \frac{1}{\rho^2} \frac{\partial^2 \psi}{\partial \phi^2} + \frac{\partial^2 \psi}{\partial z^2}.
        \]
        Substituting this into the Helmholtz equation:
        \[
        \frac{1}{\rho} \frac{\partial}{\partial \rho} \left( \rho \frac{\partial \psi}{\partial \rho} \right) + \frac{1}{\rho^2} \frac{\partial^2 \psi}{\partial \phi^2} + \frac{\partial^2 \psi}{\partial z^2} + k^2 \psi = 0.
        \]
        When \(k^2\) is generalized to:
        \[
        k^2 + f(\rho) + \frac{1}{\rho^2} g(\phi) + h(z),
        \]
        the equation becomes:
        \[
        \frac{1}{\rho} \frac{\partial}{\partial \rho} \left( \rho \frac{\partial \psi}{\partial \rho} \right) + \frac{1}{\rho^2} \frac{\partial^2 \psi}{\partial \phi^2} + \frac{\partial^2 \psi}{\partial z^2} + \left(k^2 + f(\rho) + \frac{1}{\rho^2} g(\phi) + h(z) \right) \psi = 0.
        \]
        Assume the solution \(\psi(\rho, \phi, z)\) is separable, i.e.,
        \[
        \psi(\rho, \phi, z) = R(\rho)\Phi(\phi)Z(z),
        \]
        where \(R(\rho)\) depends only on \(\rho\), \(\Phi(\phi)\) depends only on \(\phi\), and \(Z(z)\) depends only on \(z\).
        Substitute \(\psi = R(\rho)\Phi(\phi)Z(z)\) into the equation:
        \[
        \frac{1}{\rho} \frac{\partial}{\partial \rho} \left( \rho \frac{\partial (R\Phi Z)}{\partial \rho} \right) + \frac{1}{\rho^2} \frac{\partial^2 (R\Phi Z)}{\partial \phi^2} + \frac{\partial^2 (R\Phi Z)}{\partial z^2} + \left(k^2 + f(\rho) + \frac{1}{\rho^2} g(\phi) + h(z)\right) R\Phi Z = 0.
        \]
        Radial Component:
        \[
        \frac{1}{\rho} \frac{\partial}{\partial \rho} \left( \rho \frac{\partial (R\Phi Z)}{\partial \rho} \right) = \frac{\Phi Z}{\rho} \frac{\partial}{\partial \rho} \left( \rho \frac{\partial R}{\partial \rho} \right) = \Phi Z \frac{1}{\rho} \frac{d}{d \rho} \left( \rho \frac{dR}{d\rho} \right).
        \]
        Angular Component:
        \[
        \frac{1}{\rho^2} \frac{\partial^2 (R\Phi Z)}{\partial \phi^2} = \frac{RZ}{\rho^2} \frac{\partial^2 \Phi}{\partial \phi^2}.
        \]
        Axial Component:
        \[
        \frac{\partial^2 (R\Phi Z)}{\partial z^2} = R\Phi \frac{d^2 Z}{dz^2}.
        \]
        Substitute these expressions back into the equation:
        \[
        \Phi Z \frac{1}{\rho} \frac{d}{d \rho} \left( \rho \frac{dR}{d\rho} \right) + \frac{RZ}{\rho^2} \frac{d^2 \Phi}{d \phi^2} + R\Phi \frac{d^2 Z}{dz^2} + \left(k^2 + f(\rho) + \frac{1}{\rho^2} g(\phi) + h(z)\right) R\Phi Z = 0.
        \]
        Divide through by \(R(\rho)\Phi(\phi)Z(z)\):
        \[
        \frac{1}{R}\frac{1}{\rho} \frac{d}{d \rho} \left( \rho \frac{dR}{d\rho} \right) + f(\rho) + \frac{1}{\Phi} \frac{1}{\rho^2} \frac{d^2 \Phi}{d \phi^2} + \frac{1}{\rho^2} g(\phi) + \frac{1}{Z} \frac{d^2 Z}{dz^2} + h(z) + k^2 = 0.
        \]
        Group terms that depend only on \(\rho\), \(\phi\), and \(z\). The terms must independently sum to constants. Write the equation as:
        \[
        \frac{1}{R}\frac{1}{\rho} \frac{d}{d \rho} \left( \rho \frac{dR}{d\rho} \right) + f(\rho) = \lambda,
        \]
        \[
        \frac{1}{\Phi} \frac{1}{\rho^2} \frac{d^2 \Phi}{d \phi^2} + \frac{1}{\rho^2} g(\phi) = \mu,
        \]
        \[
        \frac{1}{Z} \frac{d^2 Z}{dz^2} + h(z) + k^2 = -(\lambda + \mu).
        \]
        The Helmholtz equation is separable because the terms depending on \(\rho\), \(\phi\), and \(z\) can be written as independent ordinary differential equations:
        \begin{enumerate}
            \item Radial Equation
            \[
            \frac{1}{R}\frac{1}{\rho} \frac{d}{d \rho} \left( \rho \frac{dR}{d\rho} \right) + f(\rho) = \lambda.
            \]
        
            \item Angular Equation
            \[
            \frac{1}{\Phi} \frac{1}{\rho^2} \frac{d^2 \Phi}{d \phi^2} + \frac{1}{\rho^2} g(\phi) = \mu.
            \]
        
            \item Axial Equation
            \[
            \frac{1}{Z} \frac{d^2 Z}{dz^2} + h(z) + k^2 = -(\lambda + \mu).
            \]
        \end{enumerate}
        Each equation depends on a single variable, confirming that the equation is separable.
        The Helmholtz equation:
        \[
        \nabla^2 \psi + \left(k^2 + f(\rho) + \frac{1}{\rho^2} g(\phi) + h(z)\right)\psi = 0,
        \]
        is separable in circular cylindrical coordinates.

\section*{Textbook question 9.5.1}
    Verify that the following are solutions of Laplace's equation:
    \begin{enumerate}
        \item [(a)] $\psi_1 = \frac{1}{r}$ when $r \neq 0$
        \item [(b)] $\psi_2 = \frac{1}{2r} ln\left( \frac{r + z}{r - z} \right)$
    \end{enumerate}

    \subsection*{Solution, part a:}
        The Laplacian in spherical coordinates \((r, \theta, \phi)\) is given by:
        \[
        \nabla^2 \psi = \frac{1}{r^2} \frac{\partial}{\partial r} \left( r^2 \frac{\partial \psi}{\partial r} \right) + \frac{1}{r^2 \sin \theta} \frac{\partial}{\partial \theta} \left( \sin \theta \frac{\partial \psi}{\partial \theta} \right) + \frac{1}{r^2 \sin^2 \theta} \frac{\partial^2 \psi}{\partial \phi^2}.
        \]
        For axisymmetric problems (no dependence on \(\phi\)), this reduces to:
        \[
        \nabla^2 \psi = \frac{1}{r^2} \frac{\partial}{\partial r} \left( r^2 \frac{\partial \psi}{\partial r} \right) + \frac{1}{r^2 \sin \theta} \frac{\partial}{\partial \theta} \left( \sin \theta \frac{\partial \psi}{\partial \theta} \right).
        \]
        Verifying \( \psi_1 = \frac{1}{r} \)
        Compute \( \frac{\partial \psi_1}{\partial r} \)
        \[
        \frac{\partial \psi_1}{\partial r} = \frac{\partial}{\partial r} \left( \frac{1}{r} \right) = -\frac{1}{r^2}.
        \]
        Compute \( r^2 \frac{\partial \psi_1}{\partial r} \)
        \[
        r^2 \frac{\partial \psi_1}{\partial r} = r^2 \left( -\frac{1}{r^2} \right) = -1.
        \]
        Compute \( \frac{\partial}{\partial r} \left( r^2 \frac{\partial \psi_1}{\partial r} \right) \)
        \[
        \frac{\partial}{\partial r} \left( r^2 \frac{\partial \psi_1}{\partial r} \right) = \frac{\partial}{\partial r} \left( -1 \right) = 0.
        \]
        Check \(\theta\)-dependence
        Since \( \psi_1 \) does not depend on \(\theta\):
        \[
        \frac{\partial}{\partial \theta} \left( \sin \theta \frac{\partial \psi_1}{\partial \theta} \right) = 0.
        \]
        \[
        \nabla^2 \psi_1 = 0.
        \]
        Thus, \( \psi_1 = \frac{1}{r} \) satisfies Laplace’s equation.
    
    \subsection*{Solution, part b:}
        Express \( z \) in terms of spherical coordinates.
        We know \( z = r \cos \theta \). Substitute this into \( \psi_2 \):
        \[
        \psi_2 = \frac{1}{2r} \ln\left(\frac{r + r \cos \theta}{r - r \cos \theta}\right).
        \]
        Factor \( r \) from the numerator and denominator:
        \[
        \psi_2 = \frac{1}{2r} \ln\left(\frac{1 + \cos \theta}{1 - \cos \theta}\right).
        \]
        Compute \( \frac{\partial \psi_2}{\partial r} \).
        Since \( \psi_2 \) does not explicitly depend on \( r \) in the logarithmic term:
        \[
        \frac{\partial \psi_2}{\partial r} = -\frac{1}{2r^2} \ln\left(\frac{1 + \cos \theta}{1 - \cos \theta}\right).
        \]
        Compute \( r^2 \frac{\partial \psi_2}{\partial r} \)
        \[
        r^2 \frac{\partial \psi_2}{\partial r} = r^2 \left( -\frac{1}{2r^2} \ln\left(\frac{1 + \cos \theta}{1 - \cos \theta}\right) \right) = -\frac{1}{2} \ln\left(\frac{1 + \cos \theta}{1 - \cos \theta}\right).
        \]
        Compute \( \frac{\partial}{\partial r} \left( r^2 \frac{\partial \psi_2}{\partial r} \right) \)
        \[
        \frac{\partial}{\partial r} \left( r^2 \frac{\partial \psi_2}{\partial r} \right) = \frac{\partial}{\partial r} \left( -\frac{1}{2} \ln\left(\frac{1 + \cos \theta}{1 - \cos \theta}\right) \right) = 0.
        \]
        Compute the \(\theta\)-terms
        \[
        \Phi(\theta) = \ln\left(\frac{1 + \cos \theta}{1 - \cos \theta}\right).
        \]
        \(\Phi(\theta)\) depends only on \(\cos \theta\), and its second derivatives vanish in the Laplacian operator.
        \[
        \frac{\partial}{\partial \theta} \left( \sin \theta \frac{\partial \psi_2}{\partial \theta} \right) = 0.
        \]
        \[
        \nabla^2 \psi_2 = 0.
        \]
        Thus, \( \psi_2 = \frac{1}{2r} \ln\left(\frac{r+z}{r-z}\right) \) satisfies Laplace’s equation.
        Both \( \psi_1 = \frac{1}{r} \) and \( \psi_2 = \frac{1}{2r} \ln\left(\frac{r+z}{r-z}\right) \) are solutions of Laplace’s equation.

\section*{Textbook Problem 9.6.1}
    Solve the wave equation, Eq. \ref{}, subject to the indicated conditions.
    Determine $\psi(x,t)$ given that at $t=0$ $\psi_0(x) = sin(x)$ and $\frac{\partial \psi(x)}{\partial t} = cos(x)$

    \begin{equation} \label{eq: 9.89}
        \frac{1}{c^2} \frac{\partial^2 \psi}{\partial t^2} = \frac{\partial \psi}{\partial x^2}
    \end{equation}
    \subsection*{Solution:}
        The general solution of the wave equation is:
        \[
        \psi(x, t) = f(x - ct) + g(x + ct),
        \]
        where \(f(x - ct)\) and \(g(x + ct)\) represent waves traveling in opposite directions at speed \(c\).
        Apply initial conditions (\(\psi(x, 0) = \sin x\)) and substitute \(t = 0\) into the general solution:
        \[
        \psi(x, 0) = f(x) + g(x).
        \]
        Using \(\psi(x, 0) = \sin x\), we get:
        \[
        f(x) + g(x) = \sin x. \tag{1}
        \]
        Apply initial conditions (\(\frac{\partial \psi}{\partial t}(x, 0) = \cos x\))
        First, compute \(\frac{\partial \psi}{\partial t}\):
        \[
        \frac{\partial \psi}{\partial t} = -c f'(x - ct) + c g'(x + ct).
        \]
        At \(t = 0\), this becomes:
        \[
        \frac{\partial \psi}{\partial t}(x, 0) = -c f'(x) + c g'(x).
        \]
        Using \(\frac{\partial \psi}{\partial t}(x, 0) = \cos x\), we get:
        \[
        -c f'(x) + c g'(x) = \cos x.
        \]
        Solve for \(f(x)\) and \(g(x)\)
        \[
        g(x) = \sin x - f(x).
        \]
        Substitute \(g(x) = \sin x - f(x)\) into equation
        \[
        -c f'(x) + c (\sin x - f(x))' = \cos x.
        \]
        Expand the derivative of \((\sin x - f(x))\):
        \[
        -c f'(x) + c (\cos x - f'(x)) = \cos x.
        \]
        Simplify:
        \[
        -c f'(x) + c \cos x - c f'(x) = \cos x.
        \]
        Combine terms:
        \[
        -2c f'(x) + c \cos x = \cos x.
        \]
        Divide through by \(c\):
        \[
        -2 f'(x) + \cos x = \frac{\cos x}{c}.
        \]
        Simplify further:
        \[
        f'(x) = \frac{\cos x}{2}.
        \]
        Integrate \(f'(x)\):
        \[
        f(x) = \frac{1}{2} \sin x + C_1.
        \]
        Substitute \(f(x)\) into equation:
        \[
        g(x) = \sin x - \left(\frac{1}{2} \sin x + C_1\right).
        \]
        Simplify:
        \[
        g(x) = \frac{1}{2} \sin x - C_1.
        \]
        Construct the solution using the general solution \(\psi(x, t) = f(x - ct) + g(x + ct)\), substitute \(f(x)\) and \(g(x)\) from equations:
        \[
        \psi(x, t) = \left(\frac{1}{2} \sin(x - ct) + C_1\right) + \left(\frac{1}{2} \sin(x + ct) - C_1\right).
        \]
        Simplify:
        \[
        \psi(x, t) = \frac{1}{2} \sin(x - ct) + \frac{1}{2} \sin(x + ct).
        \]
        Combine terms:
        \[
        \psi(x, t) = \frac{1}{2} \left[\sin(x - ct) + \sin(x + ct)\right].
        \]
        Using the trigonometric identity:
        \[
        \sin(A) + \sin(B) = 2 \sin\left(\frac{A + B}{2}\right) \cos\left(\frac{A - B}{2}\right),
        \]
        we get:
        \[
        \psi(x, t) = \sin(x) \cos(ct).
        \]
        The solution to the wave equation is:
        \[
        \boxed{\psi(x, t) = \sin(x) \cos(ct).}
        \]

\section*{Textbook Problem 9.7.3}
    Solve the PDE
    \[
        \frac{\partial \psi}{\partial t} = a^2 \frac{\partial^2 \psi}{\partial x^2}
    \]
    to obtain $\psi(x,t)$ for a rod of infinite extent (in both the $+x$ and $-x$ directions), with heat pulse at time $t=0$ that corresponds to $\psi_0(x) = A \delta(x)$

    \subsection*{Solution:}
        The Fourier transform of \(\psi(x, t)\) with respect to \(x\) is defined as:
        \[
        \hat{\psi}(k, t) = \int_{-\infty}^\infty \psi(x, t) e^{-ikx} \, dx,
        \]
        where \(k\) is the Fourier variable. The inverse transform is:
        \[
        \psi(x, t) = \frac{1}{2\pi} \int_{-\infty}^\infty \hat{\psi}(k, t) e^{ikx} \, dk.
        \]
        Taking the Fourier transform of both sides of the heat equation:
        \[
        \frac{\partial \psi}{\partial t} = a^2 \frac{\partial^2 \psi}{\partial x^2},
        \]
        we use the properties of the Fourier transform:
        \[
        \frac{\partial \psi}{\partial t} \mapsto \frac{\partial \hat{\psi}}{\partial t}, \quad \frac{\partial^2 \psi}{\partial x^2} \mapsto -k^2 \hat{\psi}.
        \]
        The transformed PDE becomes:
        \[
        \frac{\partial \hat{\psi}}{\partial t} = -a^2 k^2 \hat{\psi}.
        \]
        The transformed PDE is an ordinary differential equation (ODE) in \(t\):
        \[
        \frac{\partial \hat{\psi}}{\partial t} + a^2 k^2 \hat{\psi} = 0.
        \]
        The solution to this first-order ODE is:
        \[
        \hat{\psi}(k, t) = \hat{\psi}(k, 0) e^{-a^2 k^2 t}.
        \]
        Apply the initial condition by taking the Fourier transform of \(\psi(x, 0) = A \delta(x)\). The Fourier transform of \(\delta(x)\) is \(1\), so:
        \[
        \hat{\psi}(k, 0) = A.
        \]
        Thus, the solution in the Fourier domain becomes:
        \[
        \hat{\psi}(k, t) = A e^{-a^2 k^2 t}.
        \]
        To obtain \(\psi(x, t)\), apply the inverse Fourier transform:
        \[
        \psi(x, t) = \frac{1}{2\pi} \int_{-\infty}^\infty \hat{\psi}(k, t) e^{ikx} \, dk.
        \]
        Substitute \(\hat{\psi}(k, t) = A e^{-a^2 k^2 t}\):
        \[
        \psi(x, t) = \frac{A}{2\pi} \int_{-\infty}^\infty e^{-a^2 k^2 t} e^{ikx} \, dk.
        \]
        Combine the exponents:
        \[
        \psi(x, t) = \frac{A}{2\pi} \int_{-\infty}^\infty e^{-a^2 k^2 t + ikx} \, dk.
        \]
        Rewrite the exponent:
        \[
        -a^2 k^2 t + ikx = -a^2 t \left(k^2 - \frac{ikx}{a^2 t}\right).
        \]
        Complete the square in \(k\):
        \[
        k^2 - \frac{ikx}{a^2 t} = \left(k - \frac{ix}{2a^2 t}\right)^2 - \frac{x^2}{4a^4 t^2}.
        \]
        The integral becomes:
        \[
        \psi(x, t) = \frac{A}{2\pi} e^{-\frac{x^2}{4a^2 t}} \int_{-\infty}^\infty e^{-a^2 t \left(k - \frac{ix}{2a^2 t}\right)^2} \, dk.
        \]
        Change variables: Let \(u = k - \frac{ix}{2a^2 t}\), so \(dk = du\). The integral becomes:
        \[
        \psi(x, t) = \frac{A}{2\pi} e^{-\frac{x^2}{4a^2 t}} \int_{-\infty}^\infty e^{-a^2 t u^2} \, du.
        \]
        The Gaussian integral is:
        \[
        \int_{-\infty}^\infty e^{-a^2 t u^2} \, du = \sqrt{\frac{\pi}{a^2 t}}.
        \]
        Thus:
        \[
        \psi(x, t) = \frac{A}{2\pi} e^{-\frac{x^2}{4a^2 t}} \sqrt{\frac{\pi}{a^2 t}}.
        \]
        Simplify:
        \[
        \psi(x, t) = \frac{A}{\sqrt{4\pi a^2 t}} e^{-\frac{x^2}{4a^2 t}}.
        \]
        The solution to the heat equation is:
        \[
        \boxed{\psi(x, t) = \frac{A}{\sqrt{4\pi a^2 t}} e^{-\frac{x^2}{4a^2 t}}.}
        \]

\end{document}
